
\documentclass[12pt]{article}
\usepackage{geometry} % see geometry.pdf on how to lay out the page. There's lots.
\usepackage{amsmath}
\usepackage{amsfonts}
\usepackage{amssymb}
\usepackage{graphicx}
\usepackage{float}
\usepackage{cancel}
\usepackage{hyperref}
\usepackage{cleveref}
\usepackage{rotating}
\usepackage{multirow}
\usepackage{bm}

\usepackage{color}
\newcommand{\tkDM}[1]{\textcolor{red}{#1}}                     % Doddy


\geometry{a4paper} % or letter or a5paper or ... etc
% \geometry{landscape} % rotated page geometry
\newcommand{\beq}{\begin{equation}}
\newcommand{\eeq}{\end{equation}}
\newcommand{\beqa}{\begin{eqnarray}}
\newcommand{\eeqa}{\end{eqnarray}}
\newcommand{\bit}{\begin{itemize}}
\newcommand{\eit}{\end{itemize}}
% See the ``Article customise'' template for come common customisations

\title{Model Selection}
\author{}
\date{} % delete this line to display the current date

%%% BEGIN DOCUMENT
\begin{document}

\maketitle

\section{Masses, Decay Constants and Initial Conditions}

This is to guide us how hyper-parameters are related to physical parameters in each models. 
In general, the matrix is defined at the Lagrangian level as
\beqa
\mathcal{L} & = & {1\over 2} K_{ij} \partial_{\mu} \phi^i \partial^{\mu} \phi^j + {1\over 2} M_{ij} \phi^i\phi^j \\
& = & {1\over 2} \partial_{\mu} {\bm \phi}^T {\bf K} \partial_{\mu} {\bm \phi} + {1\over 2} {\bm \phi}^T {\bf M} {\bm \phi} 
\eeqa

Then after the diagonalisation of the Kahler metric with the rotation matrix ${\bm \phi} =  {\bf U} {\bm \phi'}$, the Lagrangian is in the form
\beqa
\mathcal{L} &=& {1\over 2} \partial_{\mu} {\bm \phi'}^T {\bf U}^T {\bm K} {\bf U} \partial_{\mu} {\bm \phi'} + {1\over 2} {\bm \phi'}^T {\bf U}^T {\bf M}  {\bf U} {\bm \phi'} \\
&=& {1\over 2} \partial_{\mu} {\bm \phi'}^T {\bf K}_D \partial_{\mu} {\bm \phi'} + {1\over 2} {\bm \phi'}^T {\bf U}^T {\bf M}  {\bf U} {\bm \phi'}
\eeqa
This is the step that we define decay constants, $f_{eff}^i \equiv f_i$, as
\beq
{\bf K}_D = \text{diag}(f_{1}^2, \ldots, f_{n}^2)
\eeq
that is the effective decay constants squared are the eigenvalues of the Kahler matrix.

Then, rescale the fields with
\beq
{\bm \phi'} = {\bf K}_{reciprocal} {\bm \phi''},\;\; {\bf K}_{reciprocal} = \text{diag}({1\over \sqrt{K^D_{11}}}, \ldots,  {1\over \sqrt{K^D_{nn}}})
\eeq
such that
\beq
{\bf K}_{reciprocal}^T {\bf K}_D {\bf K}_{reciprocal} = \bf 1
\eeq
and
\beq
{\bm \phi} =  {\bf U} {\bf K}_{reciprocal} {\bm \phi''}
\eeq

At this stage the Lagrangian becomes
\beq
\mathcal{L} = {1\over 2} \partial_{\mu} {\bm \phi''}^T \partial_{\mu} {\bm \phi''} + {1\over 2} {\bm \phi''}^T {\bf K}_{reciprocal}^T {\bf U}^T {\bf M}  {\bf U} {\bf K}_{reciprocal}{\bm \phi''}
\eeq

Before shift symmetry is broken, the canonically normalised axions are randomly distributed everywhere in field space. So we give the initial conditions for ${\bm \phi''}$ and $\partial_{\mu} {\bm \phi''}$ at this step as
\beqa
{\bm \phi''}_{in} = \begin{pmatrix} \theta_{in}^1 f_1 \\ \vdots \\ \theta_{in}^n f_n \end{pmatrix},\;\; \dot{\bm \phi''}_{in} =\begin{pmatrix} 0 \\ \vdots \\ 0 \end{pmatrix}
\eeqa
where $\theta^i_{in}$ is uniformly distributed from $-\pi$ to $\pi$
\beq
\theta^i_{in} \in [-\pi, \pi]
\eeq

Finally we apply the rotation matrix ${\bm \phi''} = {\bf V} \widetilde{\bm \phi}$ such that the final mass matrix is diagonal.
\beq
{\bf V}^T {\bf K}_{reciprocal}^T {\bf U}^T {\bf M}  {\bf U} {\bf K}_{reciprocal} {\bf V} = {\bf M}_{D}
\eeq
and the Lagrangian is nice and simple
\beq
\mathcal{L} = {1\over 2} \partial_{\mu} \widetilde{\bm \phi}^T \partial_{\mu} \widetilde{\bm \phi} + {1\over 2} \widetilde{\bm \phi}^T {\bf M}_{D} \widetilde{\bm \phi}
\eeq

\subsection{Summary}
In the basis of $ \widetilde{\bm \phi}$ the mass is can be found by
\beq
\boxed{ {\bf M}_{D} = {\bf V}^T {\bf K}_{reciprocal}^T {\bf U}^T {\bf M}  {\bf U} {\bf K}_{reciprocal} {\bf V}}
\eeq
where ${\bm U}$ and ${\bm V}$ are diagonalisation matrices of ${\bm K}$ and ${\bf K}_{reciprocal}^T {\bf U}^T {\bf M}  {\bf U} {\bf K}_{reciprocal}$ respectively and 
\beq
\boxed{ {\bf K}_{reciprocal} = \text{diag}({1\over \sqrt{K^D_{11}}}, \ldots,  {1\over \sqrt{K^D_{nn}}})}
\eeq
The effective decay constant can be found from
\beq
\boxed{ {\bf K}_D = \text{diag}(f_{1}^2, \ldots, f_{n}^2)}
\eeq
The initial conditions in the basis of $\widetilde{\bm \phi}$ is written as
\beq
\boxed{ \widetilde{\bm \phi}_{in} = {\bf V}^T {\bm \phi''}_{in} = {\bf V}^T \begin{pmatrix} \theta_{in}^1 f_1 \\ \vdots \\ \theta_{in}^n f_n\end{pmatrix}}
\eeq

\section{Models}

We need two sets of mean and variance for each model, i.e., $(a_0,\sigma_a)$ and $(b_0,\sigma_b)$. The exception is made in MT model where  the hyper parameters of K is $(a_0, a_1, a_2)$

\subsection{Easther-McAllister model (EM)}

In the code the tag name of the model is EM.

The Kahler matrix is assumed to be
\beq
K = \bf 1
\eeq

The potential matrix is assumed to be
\beqa
M = B\times B^T,\;\; B = \begin{pmatrix} r_{11} & \ldots & r_{1n} \\ \vdots & \ddots & \vdots \\ r_{n1} & \ldots & r_{nn} \end{pmatrix}
\eeqa
where $r_{ij}$ are i.i.d values drawn from the normal distribution with mean 0 and variance $\sigma_b^2$. $B$ is a $(N +P)×N$ rectangular matrix. This will give the marchenko pastur density function for the eigenvalues of $M$.

\subsection{Logflat model (LF)}

\tkDM{Parameterisation used so far is log normal NOT log flat. I have changed to log flat. Also minus sign is now correct.}

In the code the tag name of the model is LF.

The Kahler matrix is assumed to be
\beqa
K = A\times A^T,\;\; A = \begin{pmatrix} e_{11} & \ldots & e_{1n} \\ \vdots & \ddots & \vdots \\ e_{n1} & \ldots & e_{nn} \end{pmatrix}
\eeqa
where $e_{ij} = e^{r_{ij}}$ and $r_{ij}$ are uniform random number ranging from
\beq
r_{ij} \in [\ln a_0 -\sigma_a, \ln a_0 + \sigma_a]
\eeq
\tkDM{
\beq
r_{ij} \in U[k_{\rm min},k_{\rm max}]
\eeq
}

The potential matrix is assumed to be
\beqa
M = B\times B^T,\;\; B = \begin{pmatrix} e_{11} & \ldots & e_{1n} \\ \vdots & \ddots & \vdots \\ e_{n1} & \ldots & e_{nn} \end{pmatrix}
\eeqa
where $e_{ij} = e^{r_{ij}}$ and $r_{ij}$ are uniform random number ranging from
\beq
r_{ij} \in [\ln b_0 -\sigma_b, \ln b_0 + \sigma_b]
\eeq
\tkDM{
\beq
r_{ij} \in U[m_{\rm min},m_{\rm max}]
\eeq
}


\subsection{M-theory inspired model (MT)}

In the code the tag name of the model is MT.

The Kahler matrix is assumed to be
\beq
K = A\times A^T,\;\;
A = a_0  \begin{pmatrix} {1\over s_{11}} & \ldots & {1\over s_{1n}} \\ \vdots & \ddots & \vdots \\ {1\over s_{n1}} & \ldots & {1\over s_{nn}} \end{pmatrix}
\eeq
where $s_i$ are uniform random numbers ranging from
\beq
s_i \in [a_1,a_2]
\eeq

The potential matrix is assumed to be
\beqa
M = B\times B^T,\;\; B = \begin{pmatrix} e_{11} & \ldots & e_{1n} \\ \vdots & \ddots & \vdots \\ e_{n1} & \ldots & e_{nn} \end{pmatrix}
\eeqa
where $e_{ij} = e^{r_{ij}}$ and $r_{ij}$ are uniform random number ranging from
\beq
r_{ij} \in [\ln b_0 -\sigma_b, \ln b_0 + \sigma_b]
\eeq

\end{document}






